\section{Reduction to slice representations}

Let $G \acts V$ be a complex finite-dimensional representation of a finite group $G$.
Suppose that $\si=(\si_1,\dots,\si_n)$ is a system of homogeneous basic invariants.
Let $V^G = \{v \in V \colon Gv=v\}$ be the linear subspace of invariant vectors.
It is the subspace of all vectors $v$ for which the isotropy subgroup $G_v = \{g \in G \colon gv = v\}$ is
equal to $G$.

\subsection{Removing invariant vectors} %\label{ssec:fixedpoints}

Since finite groups are linearly reductive,
there exists a unique subrepresentation $V' \subseteq V$ such that $V = V^G \oplus V'$ (cf.\ \cite[Theorem 2.2.5]{DK02}).
Then
$\C[V]^G = \C\big[V^G\big] \otimes \C[V']^G$ and
$V/ G = V^G \times V' / G$.
A system of basic invariants of $\C[V]^G$ is given by a system of linear coordinates on $V^G$ together with a
system of basic invariants of $\C[V']^G$. Hence the following lemma is immediate.

\begin{Lemma} \label{fix}
Any lift $\overline f$ of a mapping $f=(f_0,f_1)$ in $V^G \times V' / G$ has the form
$\overline f = (f_0,\overline f_1)$,
where $\overline f_1$ is a lift of $f_1$.
\end{Lemma}


Consequently, we may assume without loss of generality that $V^G = \{0\}$.


\subsection{Luna's slice theorem}
Let us recall Luna's slice theorem.
Here we just assume that $V$ is a rational representation of~a~linearly reductive group $G$.
The categorical quotient $\pi \colon V \to V\cq G$ is the affine variety with~the coordinate ring $\C[V]^G$
together with the projection $\pi$ induced by the inclusion \mbox{$\C[V]^G \hookrightarrow \C[V]$}. In~this setting $\pi$ does not separate orbits, but for each element $z \in V \cq G$ there is a unique
closed orbit in the fiber $\pi^{-1}(z)$.
If $Gv$ is a closed orbit, then $G_v$ is again linearly reductive. We~say that $U \subseteq V$ is \emph{$G$-saturated} if $\pi^{-1} (\pi(U)) = U$.

\begin{Theorem}[{\cite{Luna73}, \cite[Theorem 5.3]{Schwarz80}}] \label{Luna}
Let $Gv$ be a closed orbit.
Choose a $G_v$-splitting \mbox{$T_v (Gv) \oplus N_v$} of $V \cong T_v V$ and let $\vh$ denote the mapping
\begin{gather*}
 G \times_{G_v} N_v \to V, \qquad [g,n] \mapsto g(v+n).
\end{gather*}
There is an affine open $G$-saturated subset $U$ of $V$ and an affine open $G_v$-saturated neighborhood~$B_v$ of $0$ in $N_v$
such that
\begin{gather*}
 \vh \colon\ G \times_{G_v} B_v \to U
\end{gather*}
and the induced mapping
\begin{gather*}
 \bar \vh \colon\ (G \times_{G_v} B_v)\cq G \to U\cq G
\end{gather*}
are \'etale.
Moreover, $\vh$ and the natural mapping $G \times_{G_v} B_v \to B_v \cq G_v$
induce a $G$-isomorphism of $G \times_{G_v} B_v$ with $U \times_{U\cq G} (B_v\cq G_v)$.
\end{Theorem}

\begin{Corollary}[{\cite{Luna73}, \cite[Corollary 5.4]{Schwarz80}}] \label{HST}
 In the setting of Theorem~$\ref{Luna}$,
 $G_y$ is conjugate to a subgroup of $G_v$ for all $y \in U$.	
 Choose a $G$-saturated neighborhood $\ol B_v$ of $0$ in $B_v$ $($classical topology$)$
 such that the canonical mapping $\ol B_v \cq G_v \to \ol U \cq G$ is a complex analytic isomorphism,
 where $\ol U = \pi^{-1}\big(\bar \vh((G \times_{G_v} B_v)\cq G)\big)$.
 Then $\ol U$ is a $G$-saturated neighborhood of $v$ and $\vh \colon G \times_{G_v} \ol B_v \to \ol U$
 is biholomorphic.
\end{Corollary}

\subsection{Uniform slice reduction} \label{ssec:red}

Let $\{\ta_i\}_{i=1}^m$ be a system of generators of $\C[N_v]^{G_v}$ and
let $\ta = (\ta_1,\dots,\ta_m) \colon N_v \rightarrow \mathbb{C}^m$ be the associated mapping.
Consider the \emph{slice}
\begin{gather*}%\label{eq:normalslice}
 S_v := v+\ol B_v,
\end{gather*}
where $\ol B_v$ is the neighborhood from Corollary~\ref{HST}.



\begin{Lemma} \label{lem:red}
 Let $a=(a_1,\dots,a_n)$ be a curve in $\si(V)$ with $a_k \ne 0$ and such that the curve
 \begin{gather*}
 \underline a := \big({a_k}^{-d_1/d_k} a_1,\dots,{a_k}^{-d_n/d_k} a_n\big)
 \end{gather*}
 lies in $\si(U_v)$, where $U_v$ is a neighborhood of $v$ in $S_v$.
 Composition of the curve $\underline a -\si(v)$ with the analytic isomorphism of Corollary~$\ref{HST}$ gives a
 curve $\underline b = (\underline b_1,\dots,\underline b_m)$ in $\ta(U_v-v)$ and
 \begin{gather*}
 b = (b_1,\dots,b_m)
 := \big(a_k^{e_1/d_k} \underline b_1,\dots,{a_k}^{e_m/d_k} \underline b_m\big),
 \qquad e_i = \deg \ta_i,
 \end{gather*}
 is a curve in $\ta(N_v)$. If $\overline b$ is a lift of $b$ over $\ta$ then
 \begin{gather*} %\label{eq:red}
 a_k^{1/d_k} v + \overline b
 \end{gather*}
 is a lift of $a$ over $\si$.
\end{Lemma}

\begin{proof}
 The curve $a_k^{-1/d_k} \overline b$ is a lift of $\underline b$ over $\ta$, indeed by homogeneity,
 \begin{gather*}
 \ta_i\big(a_k^{-1/d_k} \overline b\big) = a_k^{-e_i/d_k} \ta_i\big(\overline b\big)
 = a_k^{-e_i/d_k} b_i = \underline b_i.
 \end{gather*}
 Thus ${a_k}^{-1/d_k} \overline b + v$ is a lift of $\underline a$ over $\si$. By~homogeneity, we find $\si_i\big(\overline b + {a_k}^{1/d_k} v\big) = {a_k}^{d_i/d_k} \underline a_i = a_i$
 as required.
\end{proof}


The following lemma shows that the maximal degree of the basic invariants does not increase by
passing to a slice representation. It can be shown in analogy to \cite[Lemma 2.4]{KLMR06} or \cite{ParusinskiRainer14}.

\begin{Lemma} \label{lem:e}
 Assume that the systems of basic invariants $\{\si_j\}_{j=1}^n$ and $\{\ta_i\}_{i=1}^m$ are minimal and
 set $e := \max_{i} e_i = \max_{i} \deg \ta_i$. Then $e \le d$.
\end{Lemma}

In order to make the slice reduction uniform, we consider
the set
\begin{gather} \label{eq:compactK}
 K:= \bigg(\bigcup_{k=1}^n \big\{(\ul a_1, \dots,\ul a_n) \in \C^{n} \colon \ul a_k=1,\,|\ul a_j| \le 1 \text{ for }j \ne k\big\}\bigg) \cap \si(V),
\end{gather}
which is compact, since $\si(V)$ is closed. For each point $\underline p \in K$ choose $v \in \si^{-1}(\underline p)$.
Then the collection $\{\si(U_v)\}$ for all such $v$ is a cover of $K$ by sets $\si(U_v)$ that are open in the
trace topology on $\si(V)$
and on which
the conclusion of Lemma~\ref{lem:red} holds.
Choose a finite subcover
\begin{gather*} %\label{eq:red5}
 \cB := \{B_\de\}_{\de \in \De} = \{\si(U_{v_\de})\}_{\de \in \De}.
\end{gather*}
Then there exists $\rh>0$ such that for every $\underline p \in K$ there is a $\de \in \De$ such that
\begin{gather} \label{eq:redball}
 B_\rh(\underline p) \cap \si(V) \subseteq B_{\de},
\end{gather}
where $B_\rh(\underline p)$ is the open ball with radius $\rh$ centered at $\ul p$.

\begin{Definition}
We refer to this data as the \emph{uniform slice reduction} of the representation~$G \acts V$,
in particular, we call $\rh>0$ from~\eqref{eq:redball} the \emph{uniform reduction radius}.
\end{Definition}

\section[Estimates for a curve in sigma(V)]{Estimates for a curve in $\boldsymbol{\si(V)}$} \label{aestimates}

In the next three sections we discuss preparatory lemmas for the proof of Theorem~\ref{main} which is then
given in Section~\ref{proof}.

\subsection{An interpolation inequality}
For an interval $I \subseteq \R$ and a function $f \colon I \to \C$ we set
\begin{gather*}
 V_I(f) := \sup_{t,s \in I} |f(t)-f(s)| = \on{diam} f(I).
\end{gather*}

\begin{Lemma}[{\cite[Lemma 4]{ParusinskiRainer15}}] \label{taylor}
 Let $I \subseteq \R$ be a bounded open interval, $m \in \N_{>0}$, and $\al \in (0,1]$.
 If $f\in C^{m,\al}(\overline I)$, then for all $t\in I$
 and $s = 1,\dots,m$,
 \begin{gather*}
 \big|f^{(s)}(t) \big| \le C |I|^{-s} \bigl(V_I(f) + V_I(f)^{(m+\al-s)/(m+\al)} \big(\Hoeld_{\al,I}\big(f^{(m)}\big)\big)^{s/(m+\al)} |I|^s
 \bigr),
 \end{gather*}
 for a universal constant $C$ depending only on $m$ and $\al$.
\end{Lemma}

\subsection{The local setup}

Let $G \acts V$ be a complex finite-dimensional representation of a finite group $G$.
Assume $V^G = \{0\}$. Let $\si=(\si_1,\dots,\si_n)$ be a system of homogeneous basic invariants
of degrees $d_1,\dots, d_n$ and let $d := \max_j d_j$.
Let $a \in C^{d-1,1}\big(\ol I,\si(V)\big)$, where $I \subseteq \R$ is a bounded open interval.

It will be crucial to consider the radicals $a_j^{1/d_j}$ of the components $a_j$ of $a$ which is justified by the following remark.

\begin{Remark} \label{rem:notation}
Every continuous selection $f$ of
the multi-valued function $a_j^{1/d_j}$ is absolutely continuous on~$I$, by Theorem~\ref{cor:radicals}.
(Clearly, continuous selections exist in this case.)
Moreover, $\|f'\|_{L^1(I)}$ is independent of the
choice of the selection. Indeed, if $g$ is a different continuous selection then on each connected component $J$ of
$I \setminus \{t \colon a_j(t)=0\}$ the functions $f$ and $g$ just differ by multiplication with a fixed $d_j$-th root of unity.
Thus $\|f'\|_{L^1(J)} = \|g'\|_{L^1(J)}$. The~$\C$-valued version of Lemma~\ref{lem:extend} implies that $\|f'\|_{L^1(I)} = \|g'\|_{L^1(I)}$.

Henceforth we fix one continuous selection of $a_j^{1/d_j}$ and
denote it by
\begin{gather*}
 \we a_j \colon\ I \to \C
\end{gather*}
as well as, abusing notation, by $a_j^{1/d_j}$. We~will also consider the absolutely continuous curve
\begin{gather*}
 \we a = (\we a_1,\dots,\we a_n) \colon\ I \to \C^n.
\end{gather*}
\end{Remark}






Suppose that $t_0 \in I$ and $k \in \{1,\dots,n\}$ are such that
\begin{gather} \label{eq:k}
 |\we a_k(t_0)| = \max_{1 \le j \le n} |\we a_j(t_0)| \ne 0.
 \end{gather}
Assume further that, for some fixed positive constant $B < 1/3$,
\begin{gather}\label{assumption}
 \|\we a'\|_{L^1 (I)} \le B |\we a_k(t_0)|.
\end{gather}
At this point we just demand that the constant $B$ is fixed and smaller than $1/3$; in the proof of Theorem~\ref{main}
we will additionally specify $B$ depending on the uniform reduction radius $\rh>0$ and the maximal degree $d$, see~\eqref{eq:constB},
and it is going to be fixed along said proof. In~accordance with~\eqref{vectorvalued}, $\|\we a'\|_{L^1 (I)} = \sum_{j=1}^n \|\we a_j'\|_{L^1 (I)}$.

\begin{Definition}
 By \emph{admissible data} for $G \acts V$ me mean a tuple $(a, I, t_0, k)$, where
 $a \in C^{d-1,1}(\ol I,\si(V))$ is a curve in $\si(V)$ for a representation $G \acts V$ with $V^G = \{0\}$
 defined on an open bounded interval $I$ such that
 $t_0 \in I$ and $k \in \{1,\dots,n\}$ satisfy~\eqref{eq:k} and~\eqref{assumption}.
\end{Definition}




\subsection[The reduced curve a]{The reduced curve $\boldsymbol{\ul a}$}

Let $(a, I, t_0, k)$ be admissible data for $G \acts V$. We~shall see in the next lemma that
$a_k$ does not vanish on the interval $I$ and so the curve
\begin{align}
 \underline a \colon\ I &\to \{(\underline a_1,\dots,\underline a_n) \in \C^{n} \colon \underline a_k=1\}, \notag
 \\
 t &\mapsto \ul a(t) := \big( a_k^{-d_1/d_k} a_1, \dots, a_k^{-d_n/d_k} a_n\big)(t)
 = \big( \big(\we a_k^{-1} \we a_1\big)^{d_1}, \dots, (\we a_k^{-1} \we a_n)^{d_n}\big)(t)
 \label{curve}
\end{align}
is well-defined. The homogeneity of the basic invariants implies that $\ul a(I) \subseteq \si(V)$.

\begin{Lemma} \label{lem1}
 Let $(a, I, t_0, k)$ be admissible data for $G \acts V$.
 Then for all $t \in I$ and $j=1,\dots,n$,
 \begin{gather} \label{eq:ass10}
 |\we a_j (t) - \we a_j (t_0)| \le B |\we a_k (t_0)|,
 \\[1ex]
 \frac 2 3 < 1-B \le \bigg | \frac{\we a_k(t)}{\we a_k(t_0)} \bigg | \le 1+B < \frac4 3,\label{eq:ass11}
 \\[1ex]
 |\we a_j(t)| \le \frac 4 3 |\we a_k(t_0)| \le 2 |\we a_k(t)|.\label{eq:ass12}
 \end{gather}
 The length of the curve $\ul a$
 is bounded by $3 d^2\, 2^{d} B$.
\end{Lemma}

\begin{proof}
 First~\eqref{eq:ass10} is a consequence of~\eqref{assumption},
 \begin{gather*}
 \big|\we a_j (t) - \we a_j (t_0)\big| = \bigg|\int_{t_0}^t \we a_j' \,{\rm d}s\bigg|
 \le \| \we a_j'\|_{L^1(I)} \le B |\we a_k(t_0)|.
 \end{gather*}
 Setting $j=k$ in~\eqref{eq:ass10} easily implies~\eqref{eq:ass11}. Together with~\eqref{eq:k}, the inequalities~\eqref{eq:ass10} and~\eqref{eq:ass11} give~\eqref{eq:ass12}. In~order to estimate the length of $\ul a$ observe that
 \begin{gather*}
 \ul a_j' = \p_t \big(\big(\we a_k^{-1} \we a_j\big)^{d_j} \big) = d_j \big(\we a_k^{-1} \we a_j\big)^{d_j-1} \big(\we a_k^{-1} \we a_j'
 - \we a_k^{-2} \we a_j \we a_k' \big).
 \end{gather*}
 Since $|\we a_k^{-1} \we a_j| \le 2$, by~\eqref{eq:ass12}, and thanks to~\eqref{eq:ass11} we obtain
 \begin{gather*}
 |\ul a_j'| \le 3 d\, 2^{d} |\we a_k(t_0)|^{-1} \big( |\we a_j'| + |\we a_k'|\big).
 \end{gather*}
 Consequently, using~\eqref{assumption},
 \begin{gather*}
 \int_I |\ul a'| \,{\rm d}s \le 3 d^2\, 2^{d} B,
 \end{gather*}
 as required.
\end{proof}





%-----------------------------------------------------------------------------------------------------------------------
\section{The estimates after reduction to a slice representation} %\label{bestimates}
%-----------------------------------------------------------------------------------------------------------------------



%-----------------------------------------------------------------------------------------------------------------------
\subsection{The reduced local setup}
%-----------------------------------------------------------------------------------------------------------------------

Let $(a, I, t_0, k)$ be admissible data for $G \acts V$ such that
for all
$j = 1,\dots,n$ and $ s = 1,\dots,d-1$,
 \begin{gather}%
 \big\|a_j^{(s)} \big\|_{L^\infty(I)} \le C(d) |I|^{-s} |\we a_k (t_0)|^{d_j},\nonumber
 \\
 \Lip_I \big(a_j^{(d-1)}\big) \le C(d) |I|^{-d} |\we a_k (t_0)|^{d_j}.
\label{est:a}
 \end{gather}
Here $C(d)$ is a positive constant which depends only on $d$; in
Lemma~\ref{lem:assThm1implyassProp3} we will see that the assumptions of Theorem~\ref{main}
imply~\eqref{est:a} on suitable intervals $I$, and the proof of Lemma~\ref{lem:assThm1implyassProp3}
will provide a specific value for $C(d)$.

Additionally,
we suppose that the curve $\underline a$ (defined in~\eqref{curve}) lies entirely in one of the balls~$B_\rh(\underline p)$ from
\eqref{eq:redball}. By~Lemma~\ref{lem:red}, we obtain a curve $b \in C^{d-1,1}\big(\ol I,\ta(W)\big)$, where $H \acts W$ with $H=G_v$ and $W=N_v$ is a
slice representation of $G \acts V$ and
\begin{gather} \label{eq:b_i}
 b_i = a_k^{e_i/d_k} \ps_i \big(a_k^{-d_1/d_k} a_1, \dots, a_k^{-d_n/d_k} a_n\big), \qquad i = 1,\dots,m,
\end{gather}
 where $e_i = \deg \ta_i$ and the $\ps_i$ are analytic functions which are bounded on their domain together with all
 their partial derivatives (this may be achieved by slightly shrinking the domain).

In accordance with Remark~\ref{rem:notation} we denote by
\begin{gather*}
 \we b_i \colon\ I \to \C
\end{gather*}
a fixed continuous selection of $b_i^{1/e_i}$. Sometimes it will also be convenient to use just the
symbol $b_i^{1/e_i}$ for $\we b_i$. We~set
\begin{gather*}
 \we b = \big(\we b_1,\dots,\we b_m\big) \colon\ I \to \C^m.
\end{gather*}
Hence~\eqref{eq:b_i} can also be written as
\begin{gather*}
 b_i = \we a_k^{e_i} \ps_i \big(\we a_k^{-d_1} a_1, \dots, \we a_k^{-d_n} a_n\big) = \we a_k^{e_i}\cdot \ps_i \o \ul a.
\end{gather*}

Thanks to Lemma~\ref{fix} we may assume that $W^H = \{0\}$.


\begin{Definition}
 By \emph{reduced admissible data} for $G \acts V$ me mean a tuple $(a, I, t_0, k; b)$, where
 $(a, I, t_0, k)$ is admissible data for $G \acts V$
 satisfying~\eqref{est:a}
 such that $\ul a$
 lies entirely in one of the balls~$B_\rh(\underline p)$ from~\eqref{eq:redball}
 and $b \in C^{d-1,1}(\ol I,\ta(W))$ is a curve resulting from Lemma~\ref{lem:red}
 and thus satisfies~\eqref{eq:b_i}.
\end{Definition}


The goal of this section is to show that the bounds~\eqref{est:a} are inherited by the curve $b$ on suitable subintervals.
This requires some preparation.

\subsection[Pointwise estimates for the derivatives of b on I]
{Pointwise estimates for the derivatives of $\boldsymbol b$ on~$\boldsymbol I$}
\begin{Lemma} %\label{lem:B}
 Let $(a, I, t_0, k; b)$ be reduced admissible data for $G \acts V$.
 Then for all
 $i=1,\dots,m$ and $s = 1,\dots,d-1$,
 \begin{gather}
 \big\| b_i^{(s)} \big\|_{L^\infty(I)} \le C |I|^{-s} |\we a_k (t_0)|^{e_i},\nonumber
 \\
 \Lip_I\big( b_i^{(d-1)}\big) \le C |I|^{-d} |\we a_k (t_0)|^{e_i},
\label{eq:b_ider}
 \end{gather}
 where $C$ is a constant depending only on $d$ and on the functions $\ps_i$.
\end{Lemma}

\begin{proof}
 Let us prove the first estimate in~\eqref{eq:b_ider}.
 Let $F$ be any $C^d$-function defined on an open set $U\subseteq \C^{n}$ that contains
 $\underline a(I)$ and assume $\|F\|_{C^d(\ol U)} < \infty$. We~claim that, for $s = 1,\dots,d-1$,
 \begin{gather} \label{eq:B2}
 \|\p_t^s (F \o \ul a)\|_{L^\infty(I)} \le C |I|^{-s},
 \end{gather}
 where
 $C$ is a constant depending only on $d$ and $\|F\|_{C^d(\ol U)}$.
 For any real exponent $r$, Fa\`a di Bruno's formula implies
 \begin{gather} \label{FaadiBruno}
 \p_t^{s} \big(a_j^{r} \big) =
 \sum_{\ell \ge 1}^{s} \sum_{\ga \in \Ga(\ell,s)} c_{\ga,\ell,r}
 \, a_j^{r-\ell} a_j^{(\ga_1)} \cdots a_j^{(\ga_\ell)},
 \end{gather}
 where $\Ga(\ell,s) = \{\ga \in \N_{>0}^\ell \colon |\ga| = s\}$ and
 \begin{gather*}
 c_{\ga,\ell,r}= \frac{s!}{\ell! \ga!} r(r-1) \cdots (r-\ell+1).
 \end{gather*}
 By~\eqref{est:a} and~\eqref{eq:ass11}, this implies for $j=k$
 \begin{align}
 \|\p_t^{s} \big( a_k^{r} \big)\|_{L^\infty(I)}
 &\le
 \sum_{\ell \ge 1}^{s} \sum_{\ga \in \Ga(\ell,s)} c_{\ga,\ell,r}
 \, \|a_k^{r-\ell}\|_{L^\infty(I)} \big\|a_k^{(\ga_1)}\big\|_{L^\infty(I)}
 \cdots \big\|a_k^{(\ga_\ell)}\big\|_{L^\infty(I)}
 \notag \\
 &\le
 C(d) \sum_{\ell \ge 1}^{s} \sum_{\ga \in \Ga(\ell,s)} c_{\ga,\ell,r}
 \, |a_k(t_0)|^{r-\ell}
 |I|^{-s} |a_k(t_0)|^{\ell}
 \notag \\
 &\le
 C(d)
 |I|^{-s} |a_k(t_0)|^{r }.\label{eq:derpower}
 \end{align}
 Together with the Leibniz formula,
 \begin{gather*}
 \p_t^s \big( a_k^{- d_j/d_k} a_j\big)
 = \sum_{q=0}^{s} \binom{s}{q} a_j^{(q)} \p_t^{s-q} \big( a_k^{- d_j/d_k} \big),
 \end{gather*}
~\eqref{eq:derpower} and~\eqref{est:a} lead to
 \begin{gather} \label{eq:underlinea}
\big \|\p_t^s \big(a_k^{- d_j/d_k} a_j\big) \big\|_{L^\infty(I)}
 \le C(d) |I|^{-s}.
 \end{gather}
 Again by the Leibniz formula,
 \begin{align*}
 \p_t (F \o \ul a)
 &= \sum_{j=1}^n ((\p_{j} F)\o \ul a)\, \p_t \big( a_k^{- d_j/d_k} a_j\big),
 \\
 \p_t^s (F \o \ul a)
 &= \sum_{j=1}^n \p_t^{s-1}\Big(((\p_{j} F)\o \ul a)\, \p_t \big( a_k^{- d_j/d_k} a_j\big)\Big)
 \\
 &= \sum_{j=1}^n \sum_{p=0}^{s-1}
 \binom{s-1}{p} \p_t^{p}((\p_{j} F)\o \ul a)\, \p_t^{s- p} \big( a_k^{- d_j/d_k} a_j\big).
 \end{align*}
 For $s=1$ we immediately get~\eqref{eq:B2}. For $1< s \le d-1$, we may argue by induction on $s$. By~induction hypothesis,
 \begin{gather*}
 \|\p_t^{p}((\p_{j} F)\o \ul a)\|_{L^\infty(I)} \le C\big(d,\|\p_{j} F\|_{C^{s}(\overline U)}\big) |I|^{-p},
 \end{gather*}
 for $p = 1,\dots,s-1$. Together with~\eqref{eq:underlinea} this entails~\eqref{eq:B2}.





 Now the first part of~\eqref{eq:b_ider} is a consequence of~\eqref{eq:b_i},~\eqref{eq:derpower} (for $r = e_i/d_k$),
 and~\eqref{eq:B2} (applied to $F= \ps_i$).




For the second part of~\eqref{eq:b_ider}
observe that for functions $f_1,\dots,f_m$ on~$I$ we have
\begin{gather*}
 \Lip_I(f_1 f_2 \cdots f_m) \le \sum_{i=1}^m \Lip_I(f_i) \|f_1\|_{L^{\infty}(I)} \cdots \widehat{\|f_i\|_{L^{\infty}(I)}}
 \cdots \|f_m\|_{L^{\infty}(I)}.
\end{gather*}
Applying it to~\eqref{FaadiBruno} and using
\begin{gather*}
 \Lip_I\big(a_j^{r-\ell}\big) \le |r-\ell| \|a_j^{r-\ell-1} \|_{L^\infty(I)} \| a_j'\|_{L^\infty(I)}
\end{gather*}
we find, as in the derivation of~\eqref{eq:derpower},
\begin{gather*}
 \Lip_I\big(\p_t^{d-1}( a_k^r)\big) \le C(d,r) |I|^{-d} |a_k(t_0)|^{r}.
\end{gather*}
As above this leads to
\begin{gather*}
 \Lip_I\big(\p_t^{d-1} \big( a_k^{- d_j/d_k} a_j\big) \big)
 \le C(d) |I|^{-d},
 \end{gather*}
and
\begin{gather*}
 \Lip_I \big(\p_t^{d-1}(F \o \ul a)\big) \le C\big(d, \|F\|_{C^d(\overline U)}\big) |I|^{-d},
\end{gather*}
and finally to the second part of~\eqref{eq:b_ider}.
\end{proof}





\subsection[Integral bounds for hat b']
{Integral bounds for $\boldsymbol{\we b'}$}


Recall that $e = \max_i e_i = \max_i \deg \ta_i$ and $e' = e/(e-1)$.

\begin{Corollary} \label{lem:Lpbak}
 Let $(a, I, t_0, k; b)$ be reduced admissible data for $G \acts V$.
 Then, for all $1 \le p < e'$ and all $i=1,\dots,m$,
 \begin{gather} \label{eq:Lpbak}
\big \|\we b_i'\big\|^*_{L^p(I)} \le C |I|^{-1} |\we a_k(t_0)|,
 \end{gather}
 for a constant $C$ which depends only on $d$, $p$, and the constant in~\eqref{eq:b_ider}.
\end{Corollary}

\begin{proof}
 Notice that, by Lemma~\ref{lem:e}, we have $e \le d$. By~\eqref{est} and~\eqref{eq:b_ider},
 \begin{align*}
 \big\| \we b_i'\big\|_{e_i',w,I}&= \big\| \big(b_i^{1/e_i}\big)'\big\|_{e_i',w,I}\le C(e_i) \max \Big\{\big(\on{Lip}_{I}\big( b_i^{(e_i-1)}\big)\big)^{1/e_i}|I|^{1/e_i'},
 \|b_i'\|_{L^\infty(I)}^{1/e_i}\Big\}
 \\
 &\le C |I|^{-1 +1/e_i'} |\we a_k (t_0)|,
\end{align*}
or equivalently,
\begin{gather*}
 \big\|\we b_i'\big\|^*_{e_i',w,I} \le C |I|^{-1} | \we a_k (t_0)|.
\end{gather*}
This entails~\eqref{eq:Lpbak} in view of~\eqref{eq:qp}.
\end{proof}


%-----------------------------------------------------------------------------------------------------------------------
\subsection[Special subintervals of I and estimates on them]
{Special subintervals of $\boldsymbol I$ and estimates on them}\label{subintervals}
%-----------------------------------------------------------------------------------------------------------------------

Let $(a, I, t_0, k; b)$ be reduced admissible data for $G \acts V$.

Suppose that $t_1 \in I$ and $\ell \in \{1,\dots,m\}$ are such that
\begin{gather} \label{eq:ell}
\big |\we b_\ell(t_1)\big| = \max_{1 \le i \le m}\big |\we b_i(t_1)\big| \ne 0.
\end{gather}
By~\eqref{eq:ass12} and~\eqref{eq:b_i}, for all $t \in I$ and $i = 1,\dots,m$,
 \begin{gather} \label{eq:b_ibound}
\big|\we b_i(t) \big| \le C_1 |\we a_k (t_0)|,
 \end{gather}
where the constant $C_1$ depends only on the functions $\ps_i$.
Thanks to~\eqref{eq:b_ibound} we can choose
a~con\-stant $D< 1/3$ and an open interval $J$ with $t_1 \in J \subseteq I$ such that
\begin{gather} \label{assumption2a}
 | J| |I|^{-1} {|\we a_k(t_0)|}
 + \big\|\we b '\big\|_{L^1 (J)} = D \big|\we b_\ell(t_1)\big|,
\end{gather}
where $\big\|\we b '\big\|_{L^1 (J)} = \sum_{i=1}^m\big \|\we b_i '\big\|_{L^1 (J)}$.
It suffices to take $D < C_1^{-1}$, where $C_1$ is the constant in~\eqref{eq:b_ibound}.
Here we use that $\we b _i$ is absolutely continuous,
by Theorem~\ref{cor:radicals}.

We will now see that on the interval $J$ the estimates of Section \ref{aestimates} hold for $b_i$
instead of $a_j$.


\begin{Lemma} \label{lemB}
 Let $(a, I, t_0, k; b)$ be reduced admissible data for $G \acts V$.
 Assume that $t_1 \in I$ and $\ell \in \{1,\dots,m\}$ are such that
~\eqref{eq:ell} holds and
 let $D$ and $J$ be as in~\eqref{assumption2a}.
 Then,
 for all $t \in J$ and $i = 1,\dots,m$,
 \begin{gather}
 \big| \we b_i (t) - \we b_i (t_1)\big| \le D \big|\we b_\ell (t_1)\big|, \label{b1}
 \\
 \frac 2 3 < 1-D \le \left | \frac{\we b_\ell(t)}{\we b_\ell(t_1)} \right| \le 1+D < \frac 4 3,
 \label{b2}
 \\
 \big|\we b_i(t)\big|\le \frac 4 3 \big|\we b_\ell(t_1)\big| \le 2 \big|\we b_\ell(t)\big|. \label{b3}
 \end{gather}
 The length of the curve
 \begin{gather*}
 J \ni t \mapsto \ul b(t)
 := \big( b_\ell^{-e_1/e_\ell} b_1, \dots, b_\ell^{-e_m/e_\ell} b_{m}\big)(t)
 = \big( \big(\we b_\ell^{-1} \we b_1\big)^{e_1}, \dots, \big(\we b_\ell^{-1} \we b_m\big)^{e_m}\big)(t)
 \end{gather*}
 in $\ta(W)$ is bounded by $3 e^2\, 2^{e} D$.
 For all $ i = 1,\dots,m$ and $ s = 1,\dots,d-1$,
 \begin{gather}
 \big\| b_i^{(s)} \big\|_{L^\infty(J)} \le C |J|^{-s} \big|\we b_\ell (t_1)\big|^{e_i},\nonumber
 \\
 \Lip_J\big(b_i^{(d-1)}\big) \le C |J|^{-d} \big|\we b_\ell (t_1)\big|^{e_i},\label{b4}
 \end{gather}
 for a universal constant $C$ depending only on $d$ and $\ps_i$.
\end{Lemma}

\begin{proof}
 The proof of~\eqref{b1}--\eqref{b3} is analogous to the proof of Lemma~\ref{lem1}; use~\eqref{eq:ell} and~\eqref{assumption2a}
 instead of~\eqref{eq:k} and~\eqref{assumption}. The bound for the length of the curve $J \ni t \mapsto \ul b(t)$ (which is
 well-defined by~\eqref{b2})
 follows from~\eqref{assumption2a}--\eqref{b3}; see the proof of Lemma~\ref{lem1}.

 Let us prove~\eqref{b4}. By~\eqref{eq:b_ider}, for $i = 1,\dots,m$ and $s= 1,\dots,d-1$
 \begin{gather}
 \big\| b_i^{(s)} \big\|_{L^\infty(I)} \le C |I|^{-s} |\we a_k (t_0)|^{e_i},\nonumber
 \\
 \Lip_I\big( b_i^{(d-1)}\big) \le C |I|^{-d} | \we a_k (t_0)|^{e_i},
 \label{eq:b_iderj}
 \end{gather}
 where $C = C(d, \ps_i)$. Recall that $e \le d$.

 For $s\ge e_i$ (including the case $s = d$), we have $\big(|J||I|^{-1}\big)^s \le \big(|J||I|^{-1}\big)^{e_i}$ and thus
 \begin{gather*}
 | I |^{-s} |\we a_k(t_0)|^{e_i}
 \le | J |^{-s} \big(|J| |I|^{-1} |\we a_k(t_0)| \big)^{e_i}
 \le | J |^{-s} \big|\we b_\ell(t_1)\big|^{e_i},
 \end{gather*}
 where the second inequality follows from~\eqref{assumption2a}. Hence~\eqref{eq:b_iderj} implies~\eqref{b4}.


 For $t\in J$ and $s< e_i$,
 \begin{align*}
 \big|b_i^{(s)}(t) \big|
 &\le C | J |^{-s} \bigl ( V_J(b_i)
 + V_J(b_i)^{(e_i-s)/e_i} \big(\Lip_J \big(b_i^{(e_i-1)}\big)\big)^{s/e_i} |J|^s \bigr)
 \hspace{7mm} \text{by Lemma~\ref{taylor}}
 \\
 &\le C_1 | J |^{-s} \Bigl ( \big|\we b_\ell(t_1)\big|^{e_i}
 + |\we b_\ell(t_1)|^{e_i-s} |J|^s |I|^{-s} |\we a_k(t_0)|^{s} \Bigr)
 \hspace{16mm} \text{by~\eqref{b3} and~\eqref{eq:b_iderj}}
 \\
 &\le C_2 | J |^{-s} \big|\we b_\ell(t_1)\big|^{e_i}, \hspace{69mm} \text{by~\eqref{assumption2a}}
 \end{align*}
 for constants $C = C(e_i)$ and $C_h = C_h(d,\ps_i)$.
\end{proof}


%-----------------------------------------------------------------------------------------------------------------------
\section{A special cover by intervals} %\label{specialcover}
%-----------------------------------------------------------------------------------------------------------------------

In the proof of Theorem~\ref{main} we shall have to glue local integral bounds on small intervals which result from the splitting
process to global bounds. In~this section we present a technical result which will
allow us to do so.

Let us suppose that $H \acts W$ is a complex finite-dimensional representation of a finite group~$H$,
$\ta= (\ta_1,\dots, \ta_m)$ is a system of homogeneous basic invariants of degree $e_i = \deg \ta_i$, and
$e := \max_i e_i$.



\subsection{Covers by prepared collections of intervals}\label{sec:preparation}
Let $I \subseteq \R$ be a bounded open interval and
let $b \in C^{e-1,1}\big(\overline I,\ta(W)\big)$.
For each point $t_1$ in
\begin{gather*}
 I' := I \setminus \{t \in I \colon b (t) = 0\}
\end{gather*}
there exists $\ell \in \{1, \dots, m\}$ such that~\eqref{eq:ell} holds.
Assume that there are positive constants $D < 1/3$ and $L$ such that
for all $t_1 \in I'$
there is
an open interval $J=J(t_1)$ with $t_1 \in J \subseteq I$ such that
\begin{gather} \label{assumption2aG}
 L | J| + \big\|\we b'\big\|_{L^1 (J)} = D\big |\we b_\ell(t_1)\big|.
\end{gather}
Note that~\eqref{eq:ell} and~\eqref{assumption2aG} imply~\eqref{b2} (cf.\ the proof of Lemma~\ref{lemB});
in particular, we have $J \subseteq I'$.

This defines a collection $\cI:=\{J(t_1)\}_{t_1\in I'}$ of open (in the relative topology) intervals which cover $I'$. We~will \emph{prepare} this collection in the following way.
Let us consider the functions
\begin{gather*}
 \vh_{t_1,+}(s) := L (s-t_1) + \big\|\we b'\big\|_{L^1 ([t_1,s))}, \qquad s\ge t_1,
 \\
 \vh_{t_1,-}(s) := L (t_1-s) + \big\|\we b'\big\|_{L^1 ((s,t_1])}, \qquad s\le t_1.
\end{gather*}
Then $\vh_{t_1,\pm} \ge 0$ are monotonic continuous functions defined for small $\pm (s - t_1)\ge 0$
and satisfying $\vh_{t_1,\pm}(t_1) = 0$.

Fix $t_1 \in I'$. Thanks to~\eqref{assumption2aG} there exist $s_-,s_+ \in \R$ such that
\begin{gather*}
\vh_{t_1,-}(s_-) + \vh_{t_1,+}(s_+) = D \big|\we b_\ell(t_1)\big|
\end{gather*}
and $J(t_1) = (s_-,s_+)$.
But there may also be a choice $s'_-,s'_+ \in \R$ such that this occurs \emph{symmetrically}, that is
\begin{gather*} %\label{1kind}
 \vh_{t_1,-}(s'_-) = \vh_{t_1,+}(s'_+) = \frac D 2\big|\we b_\ell(t_1)\big|.
 \end{gather*}
If such a choice $s'_-,s'_+ \in \R$ exists, we replace $J(t_1)$ in the collection $\cI$ by the interval $(s'_-,s'_+)$.
(In \cite{ParusinskiRainer15} we said that these are intervals of \emph{first kind}.)
If such a choice does not exist, then we leave $J(t_1)$ in $\cI$ unchanged;
this happens when
we reach the boundary of the interval $I$ before either $\vh_{t_1,-}$ or $\vh_{t_1,+}$ has
 grown to the value $(D/2)|\we b_\ell(t_1)|$. (These intervals were said to be of~\emph{second kind} in \cite{ParusinskiRainer15}.)

If a collection $\cI$ satisfies this property, we say that it is \emph{prepared}.



\subsection{A special subcollection of intervals}


\begin{Proposition} \label{cover}
 Let $I \subseteq \R$ be a bounded open interval.
 Let $b \in C^{e-1,1}(\overline I,\ta(W))$.
 For each point $t_1$ in $I'$ fix $\ell \in \{1, \dots, m\}$ such that~\eqref{eq:ell} holds.
 Let $\cI=\{J(t_1)\}_{t_1\in I'}$ be a collection of~open intervals $J=J(t_1)$ with $t_1 \in J \subseteq I'$ such that:
\begin{enumerate}\itemsep=0pt
 \item[$1.$] There are positive constants $D < 1/3$ and $L$ such that
 for all $t_1 \in I'$ we have~\eqref{assumption2aG} for $J=J(t_1)$.
 \item[$2.$] The collection $\cI$ is prepared as explained in Section~$\ref{sec:preparation}$.
 \end{enumerate}
Then the collection $\cI$ has a countable subcollection $\cJ$ that still
 covers $I'$ and such that every point in $I'$ belongs to at most two intervals in $\cJ$. In~particular,
 \begin{gather*}
 \sum_{J \in \cJ} |J| \le 2 |I'|.
 \end{gather*}
\end{Proposition}

\begin{proof}
 It follows from the proof of \cite[Proposition 2]{ParusinskiRainer15}.
\end{proof}


\begin{Remark} \label{rem:cover}
 It is essential for us that $\cJ$ is a subcollection and not a refinement; by shrinking the intervals we would
 lose equality in~\eqref{assumption2aG}. We~will need this proposition for gluing local $L^p$-estimates to global ones.
\end{Remark}





%-----------------------------------------------------------------------------------------------------------------------
\section{Proof of \texorpdfstring{Theorem~\ref{main}}{Theorem 1.1}} \label{proof}
%-----------------------------------------------------------------------------------------------------------------------

The proof is based on uniform slice reduction and induction on the order of~$G$. We~will apply the following convention:
\begin{quote}
 \it We will no longer explicitly state all the dependencies of the constants. Henceforth, their dependence on the
 data of the uniform slice reductions will be subsumed by~simply indicating that they depend on the representation
 $G \acts V$. This includes the choice of $\si$: different choices of the basic invariants yield different constants.
 The~constants which are uniform in this sense will be denoted by $C = C(G \acts V)$ and may vary from line to line.
\end{quote}



\subsection*{Outline of the proof}
The proof of Theorem~\ref{main} is divided into three steps.
\begin{enumerate}\itemsep=0pt
\setlength{\leftskip}{0.45cm}
 \item[{\it Step} 1:] We check that for any $a \in C^{d-1,1}([\al,\be],\si(V))$ and all points $t_0 \in (\al,\be)$, where \mbox{$a(t_0) \ne 0$},
 we can find $k$ and a suitable interval $I$ such that $(a|_I,I,t_0,k;b)$, where $b$ is obtained by Lemma~\ref{lem:red},
 is reduced admissible data for $G \acts V$.
 \item[{\it Step} 2:] The reduced admissible data $(a|_I,I,t_0,k;b)$ represents the hypothesis of the inductive argument which is the heart of the proof.
 It will show that every continuous lift of $b$ is~absolutely continuous on~$I$ and it will give an $L^p$-bound for the first derivative
 of the lift on~$I$.
 \item[{\it Step} 3:] We assemble the proof of Theorem~\ref{main}.
 The local bounds will be glued to global bounds for lifts of the original curve $a$.
 \end{enumerate}










%-----------------------------------------------------------------------------------------------------------------------
\subsection*{Step 1: The assumptions of Theorem~\ref{main} imply the local setup of the induction}
%-----------------------------------------------------------------------------------------------------------------------

Assume that~$V^G = \{0\}$.
Let $a \in C^{d-1,1}([\al,\be],\si(V))$. Let $\rh$ be the uniform reduction radius from~\eqref{eq:redball}. We~fix a universal positive constant $B$ satisfying
\begin{gather} \label{eq:constB}
 B < \min\bigg\{\frac1 3, \frac{\rh}{3 d^2 2^{d}}\bigg\}.
\end{gather}
Fix $t_0 \in (\al,\be)$ and $k \in \{1,\dots,n\}$ such that
\begin{gather} \label{k}
 |\we a_k(t_0)| = \max_{1 \le j \le n} |\we a_j(t_0)|\ne 0.
\end{gather}
This is possible unless $a \equiv 0$ in which case nothing is to prove.
Choose a maximal open interval $I \subseteq (\al,\be)$ containing $t_0$ such that
\begin{gather}\label{assumption<=}
\const |I| + \|\we a'\|_{L^1 (I)} \le B |\we a_k(t_0)|,
\end{gather}
where
\begin{gather} \label{M}
 \const = \max_{1 \le j\le n} \big(\Lip_I\big(a_j^{(d-1)}\big)\big)^{1/d} |\we a_k (t_0)|^{(d-d_j)/ d}.
\end{gather}

Consider the point $\underline p = \ul a(t_0)$, where $\ul a$ is the curve defined in~\eqref{curve}. By~\eqref{k},
$\underline p$ is an~ele\-ment of the set $K$ defined in~\eqref{eq:compactK}. By~the properties of the uniform slice reduction specified in~Section~\ref{ssec:red},
the ball $B_\rh(\underline p)$ is contained in some ball of the finite cover $\cB$ of $K$. By~Lemma~\ref{lem1} and~\eqref{eq:constB},
the length of the curve $\ul a|_I$ is bounded by $\rh$.
Thus
\begin{gather} \label{bspecified}
 \text{$b \in C^{d-1,1}\big(\ol I,\ta(W)\big)$ is obtained by Lemma~\ref{lem:red} and satisfies~\eqref{eq:b_i}.}
\end{gather}




\begin{Lemma} \label{lem:assThm1implyassProp3}
	Assume that $V^G = \{0\}$.
 Let $(\al,\be) \subseteq \R$ be a bounded open interval and let $a \in C^{d-1,1}([\al,\be],\si(V))$.
 Let $B$ be a positive constant satisfying~\eqref{eq:constB}.
 Let $t_0 \in (\al,\be)$ and $k \in \{1,\dots,n\}$ be such that~\eqref{k} holds.
 Let $I$ be an open interval with $t_0 \in I \subseteq (\al,\be)$ satisfying~\eqref{assumption<=}
 and $b$ the reduced curve from~\eqref{bspecified}.
 Then $(a|_I,I,t_0,k;b)$ is reduced admissible data for $G \acts V$.
\end{Lemma}

\begin{proof}
	It remains to prove~\eqref{est:a}, i.e.,
	for all
	$ j = 1,\dots,n$ and $ s = 1,\dots,d-1$,
 \begin{gather*}%\label{est:a}
 \big\|a_j^{(s)} \big\|_{L^\infty(I)} \le C\, |I|^{-s} |\we a_k (t_0)|^{d_j},
 \\
 \Lip_I \big(a_j^{(d-1)}\big) \le C\, |I|^{-d} |\we a_k (t_0)|^{d_j},
 \end{gather*}	
 for $C = C(G \acts V)$.	
 The second bound is immediate from~\eqref{assumption<=}.
Let $t \in I$. By~Lemma~\ref{taylor},
 \begin{gather*}
 \big| a_j^{(s)}(t) \big| \le C |I|^{-s} \bigl( V_I(a_j)
 + V_I(a_j)^{(d-s)/d} \Lip_I\big(a_j^{(d-1)}\big)^{s/d} |I|^s\bigr).
 \end{gather*}
By~\eqref{eq:ass12} (it is clear that $(a|_I,I,t_0,k)$ is admissible data for $G \acts V$),
\begin{gather*}
 V_I(a_j) \le 2 \|a_j\|_{L^{\infty}(I)} \le 2\, (4/3)^d |\we a_k (t_0)|^{d_j},
\end{gather*}
and, by~\eqref{assumption<=},
\begin{gather*}
 \max_{1 \le j\le n} \big(\Lip_I\big(a_j^{(d-1)}\big)\big)^{s/d} |\we a_k (t_0)|^{-d_j s/d} |I|^s
 = |\we a_k (t_0)|^{-s} \const^s |I|^s \le 1.
\end{gather*}
Thus
\begin{gather*}
\big | a_j^{(s)}(t) \big| \le C |I|^{-s}
 |\we a_k (t_0)|^{d_j} \big(C_1 + C_2 \Lip_I\big(a_j^{(d-1)}\big)^{s/d} |\we a_k (t_0)|^{-d_j s/d} |I|^s \big)
\le C_3 |I|^{-s} |\we a_k (t_0)|^{d_j},
\end{gather*}
for constants $C_i$ that depend only on $d$.
So~\eqref{est:a} is proved.
\end{proof}

\subsection*{Step 2: The inductive argument}

The heart of the proof of Theorem~\ref{main} is the following


\begin{Proposition} \label{induction}
 Let $(a, I, t_0, k; b)$ be reduced admissible data for $G \acts V$.
 Then every con\-ti\-nuous lift $\ol b \in C^0(I,W)$ of $b$ is absolutely con\-ti\-nuous and satisfies
 \begin{gather} \label{muab}
 \|\ol b'\|_{L^p(I)} \le C \Big( \| |I|^{-1} {|\we a_k(t_0)|} \|_{L^p (I)}
 + \big\|\we b'\big\|_{L^p (I)}\Big),
 \end{gather}
 for all $1 \le p < d'$ and a constant $C$ depending only on $G \acts V$ and $p$.
\end{Proposition}




\begin{Remark}
 Notice that we bound the $L^p$-norm of the derivative of a general lift $\ol b$
 by the $L^p$-norm of the derivatives of the lifts $\we b_i$ for the standard action of rotation in $\C$.
\end{Remark}



\begin{proof}[Proof of Proposition~\ref{induction}]
	We proceed by induction on the group order.


\medskip\noindent
{\bf Induction basis.}
Proposition~\ref{induction} trivially holds, if the slice representation $H \acts W$
is trivial. In~that case $\C[W]^H \cong \C[W]$ and any system $\ta = (\ta,\dots,\ta_m)$ of linear coordinates
is a minimal system of generators. Hence $\ta \colon W \to \C^m$ is a linear isomorphism.
Moreover, $e_1 = e_2 = \cdots = e_m = e =1$ whence $\we b_i = b_i$ for all $i$.
Any lift of $b \in C^{d-1,1}\big(\ol I,\ta(W)\big)$ is of the form $\ol b = \ta^{-1} \o b$ and thus
\eqref{muab} is trivially satisfied.

\medskip\noindent
{\bf Inductive step.}
Let us set
\begin{gather*}
 I' := I \setminus \{t \in I \colon b(t) = 0\}.
\end{gather*}
For each $t_1 \in I'$ choose $\ell \in \{1, \dots, m\}$ such that~\eqref{eq:ell} holds. By~Section \ref{subintervals}, there is an open interval $J = J(t_1)$, $t_1 \in J \subseteq I'$,
such that~\eqref{assumption2a}, i.e.,
\begin{gather*}
 | J| |I|^{-1} {|\we a_k(t_0)|}
 + \big\|\we b'\big\|_{L^1 (J)} = D \big|\we b_\ell(t_1)\big| .
\end{gather*}
The constant $D$
can be chosen sufficiently small such that the length of the curve $\underline b|_J$ is bounded by
the uniform reduction radius $\si$ of the representation $H \acts W$.
It suffices to take
\begin{gather} \label{D}
 D < \min\bigg\{\frac 1 3, \frac{\si}{3 e^2 2^{e}}, C_1^{-1}\bigg\},
\end{gather}
where $C_1$ is the constant in~\eqref{eq:b_ibound}.
This follows from Lemma~\ref{lemB}
and the arguments in Section~\ref{ssec:red} and in Step~1 applied to $b$.

Then Lemma~\ref{lem:red} provides a curve $c \in C^{d-1,1}\big(\ol J, \pi(X)\big)$, where $K \acts X$ is a slice representation
of $H \acts W$, $\pi = (\pi_1,\dots,\pi_q)$ is a system of homogeneous basic invariants with degrees $f_1,\dots,f_q$,
and $f = \max_h f_h$.
The components of $c$ satisfy
\begin{gather*} %\label{eq:c_h}
 c_h = b_\ell^{f_h/e_\ell} \th_h \big(b_\ell^{-e_1/e_\ell} b_1, \dots, b_\ell^{-e_m/e_\ell} b_m\big), \qquad h = 1,\dots,q,
\end{gather*}
for suitable analytic functions $\th_h$. We~adopt our usual convention that
\begin{gather*}
\we c_h \colon\ J \to \C
\end{gather*}
denotes a fixed continuous selection of $c_h ^{1/f_h}$ and
set
\begin{gather*}
	\we c = (\we c_1 ,\dots,\we c_q) \colon\ J \to \C^q.
\end{gather*}
In view of Lemma~\ref{lemB} we conclude that $(b,J,t_1,\ell; c)$ is reduced admissible data for $H \acts W$.







By Proposition~\ref{cover} (where~\eqref{assumption2a} plays the role of~\eqref{assumption2aG}),
we may conclude that there is a countable family $\{(J_\ga,t_\ga,\ell_\ga,c_\ga)\}$ of open intervals $J_\ga \subseteq I'$,
of points $t_\ga \in J_\ga$,
of integers $\ell_\ga \in \{1, \dots,m\}$,
and reduced curves $c_\ga$ such that, for all $\ga$,
\begin{itemize}
	\item $(b,J_\ga,t_\ga,\ell_\ga;c_\ga)$ is reduced admissible data for $H \acts W$,
	\item we have
	\begin{gather}
 |J_\ga| |I|^{-1} {|\hat a_k(t_0)|} + \big\|\hat b'\big\|_{L^1 (J_\ga)} =
 D \big|\hat b_{\ell_\ga}(t_\ga)\big|, \label{Jba}
\end{gather}
\item and
\begin{gather}
	\bigcup_\ga J_\ga = I', \qquad \sum_\ga |J_\ga| \le 2 |I'|. \label{Jga}
\end{gather}
\end{itemize}


Let $\ol b \in C^0(I,W)$ be a continuous lift of $b$.
Fix $\ga$ and let $K \acts X$ be the corresponding slice representation of $H \acts W$.
Since $H$ is a finite group, we have $W \cong X$. With this identification and the decomposition $X = X^K \oplus X'$
we may deduce that the component of $\ol b$ in $X'$
is a~continuous lift of~$c_\ga$ on the interval~$J_\ga$.
To simplify the notation we will assume without loss of~generality that $X^{K} = \{0\}$ and that
$\ol b$ is a lift of $c_\ga$ on the interval $J_\ga$.

The induction hypothesis implies that
$\ol b$ is absolutely continuous on $J_\ga$ and satisfies
\begin{gather} \label{mu}
\big \|\ol b'\big\|_{L^p(J_\ga)} \le C \Big( \big\| |J_\ga|^{-1} {\big|\hat b_{\ell_\ga}(t_\ga)\big|} \big\|_{L^p (J_\ga)}
 + \|\hat c_{\ga}'\|_{L^p (J_\ga)}\Big),
\end{gather}
for all $1 \le p < e'$, where $C$ is a constant depending only on $H \acts W$ and $p$.


\medskip\noindent
{\bf $\boldsymbol{L^p}$-estimates on $\boldsymbol I$.}
To finish the proof of Proposition~\ref{induction} we have to show that the estimates~\eqref{mu}
on the subintervals $J_\ga$
imply the bound~\eqref{muab} on~$I$.
To this end we observe that Corollary~\ref{lem:Lpbak} (applied to $(b,J_\ga,t_\ga,\ell_\ga;c_\ga)$)
implies that,
for all $p$ with $1 \le p < f_\ga'$,
{\samepage\begin{gather} \label{eq:ctob}
 \|\hat c_{\ga}'\|^*_{L^p(J_\ga)}
 \le C |J_\ga|^{-1} \big|\hat b_{\ell_\ga} (t_\ga)\big|,
\end{gather}
for a constant $C$ that depends only on $H \acts W$ and $p$.

}


Now~\eqref{eq:ctob} and~\eqref{Jba} allow us to estimate the right-hand side of
\eqref{mu}:
\begin{align*}
 \big\||J_\ga|^{-1} \big|\hat b_{\ell_\ga}(t_\ga)\big| \big\|^*_{L^p (J_\ga)}
 + \|\hat c_{\ga}'\|^*_{L^p (J_\ga)}
 &= |J_\ga|^{-1} \big|\hat b_{\ell_\ga}(t_\ga)\big| +
 \|\hat c_{\ga}'\|^*_{L^p (J_\ga)}
 \\
 &\le C |J_\ga|^{-1} \big|\hat b_{\ell_\ga}(t_\ga)\big|
 \\
 &= C D^{-1} \Big( \big\| |I|^{-1} {|\hat a_k(t_0)|} \big\|^*_{L^1 (J_\ga)}
 + \big\|\hat b'\big\|^*_{L^1 (J_\ga)}\Big)
 \\
 &\le C D^{-1} \Big( \big\| |I|^{-1} {|\hat a_k(t_0)|} \big\|^*_{L^p (J_\ga)}
 + \big\|\hat b'\big\|^*_{L^p (J_\ga)}\Big)
\end{align*}
and therefore
\begin{gather}
\big \||J_\ga|^{-1} \big|\hat b_{\ell_\ga}(t_\ga)\big|\big \|^p_{L^p (J_\ga)}
 + \|\hat c_{\ga}'\|^p_{L^p (J_\ga)}
 \le C D^{-p} \Big( \big\| |I|^{-1} {|\hat a_k(t_0)|} \big\|^p_{L^p (J_\ga)}
 + \big\|\hat b'\big\|^p_{L^p (J_\ga)}\Big), \label{cba}
\end{gather}
for a constant $C$ that depends only on $H \acts W$ and $p$.




Let us now glue the bounds on $J_\ga$ to a bound on~$I$. By~\eqref{Jga},~\eqref{mu}, and~\eqref{cba},
\begin{gather*} %\label{ind11}
 \sum_\ga \big\|\ol b'\big\|^p_{L^p(J_\ga)}
 \le C D^{-p} \Big( \big\| |I|^{-1} {|\hat a_k(t_0)|} \big\|^p_{L^p (I)}
 + \big\|\we b'\big\|^p_{L^p (I)}\Big),
\end{gather*}
for a constant $C$ that depends only on $H \acts W$ and $p$.
Thus $\ol b$ is absolutely continuous on $I'$
and
\begin{gather*}
 \big\|\ol b'\big\|_{L^p(I')}
 \le C D^{-1} \Big( \big\| |I|^{-1} {|\we a_k(t_0)|} \big\|_{L^p (I)}
 + \big\|\we b'\big\|_{L^p (I)}\Big),
\end{gather*}
for a constant $C$ that depends only on $H \acts W$ and $p$.
Since $\ol b$ vanishes on $I \setminus I'$,
Lemma~\ref{lem:extend} implies that $\ol b$ is absolutely continuous on~$I$ and
satisfies~\eqref{muab}, since $D = D(H \acts W)$ by~\eqref{D}.
This completes the proof of Proposition~\ref{induction}.
\end{proof}

\subsection*{Step 3: The proof of Theorem~\ref{main}}

In view of Lemma~\ref{fix} we may assume $V^G = \{0\}$.
Let $a \in C^{d-1,1}([\al,\be],\si(V))$. Suppose that $B$ is a positive constant fulfilling~\eqref{eq:constB} and
assume that
$t_0 \in (\al,\be)$, $k \in \{1,\dots,n\}$, and $I \ni t_0$
satis\-fy~\eqref{k} and~\eqref{assumption<=}.
Let $b \in C^{d-1,1}\big(\overline I, \ta(W)\big)$ be the reduced curve from~\eqref{bspecified}.
Then Lemma~\ref{lem:assThm1implyassProp3} implies that $(a,I,t_0,k;b)$ is reduced admissible data and
consequently each continuous lift $\ol b$ of~$b$ satisfies~\eqref{muab}, by Proposition~\ref{induction}. In~particular, if~$\ol a \in C^0((\al,\be),V)$ is a continuous lift of $a$, then we may assume
that $\ol a|_I$ is a lift of $b$.
It follows that $\ol a$ is absolutely continuous on~$I$ and
\begin{gather} \label{amuab}
 	\|\ol a'\|_{L^p(I)} \le C(G \acts V,p) \Big( \big\| |I|^{-1} {|\we a_k(t_0)|} \big\|_{L^p (I)}
 + \big\|\we b'\big\|_{L^p (I)}\Big).
 \end{gather}
Our next goal is to estimate the right-hand side of~\eqref{amuab} in terms of~$a$.

By Corollary~\ref{lem:Lpbak},
we get for all $p$ with $1 \le p < e'$,
\begin{gather} \label{beforecases}
 \big\||I|^{-1} |\we a_k (t_0)| \big\|^*_{L^p (I)} + \big\|\we b'\big\|^*_{L^p(I)}
 \le C |I|^{-1} |\we a_k (t_0)|,
\end{gather}
where the constant $C$ depends only on $G \acts V$ and $p$.
At this stage we distinguish the two cases of strict inequality or equality in~\eqref{assumption<=}:
\begin{enumerate}
 \item[$(i)$] Strict inequality: we have $I = (\al,\be)$ and
 \begin{gather}\label{assumption<}
 \const |I| + \|\we a'\|_{L^1 (I)} < B |\we a_k(t_0)|.
 \end{gather}
 \item[$(ii)$] Equality:
 \begin{gather}\label{assumption=}
 \const |I| + \|\we a'\|_{L^1 (I)} = B |\we a_k(t_0)|.
 \end{gather}
\end{enumerate}

\medskip\noindent
{\bf Case ($\boldsymbol i$).}
In this case we can reduce to the curve $b \in C^{d-1,1}\big(\overline I,\ta(W)\big)$ on the whole interval $I=(\al,\be)$; cf.\
Step 1. Thus,~\eqref{beforecases} becomes
\begin{gather*}
 \big\|(\be-\al)^{-1} |\we a_k (t_0)| \big\|_{L^p ((\al,\be))}
 + \big\|\we b'\big\|_{L^p((\al,\be))}
 \le C (\be-\al)^{-1+1/p} |\we a_k(t_0)|,
\end{gather*}
which can be bounded by
\begin{gather*} %\label{eq:blowup}
 C (\be-\al)^{-1+1/p} \max_{1 \le j \le n} \| a_j\|^{1/d_j}_{L^\infty((\al,\be))}.
\end{gather*}
By~\eqref{amuab}, $\ol a$ is absolutely continuous on $(\al,\be)$ and
\begin{gather} \label{lacase2}
 \|\ol a'\|_{L^p((\al,\be))} \le C (\be-\al)^{-1+1/p} \max_{1 \le j \le n} \|a_j\|^{1/d_j}_{L^\infty((\al,\be))},
\end{gather}
where $C = C(G \acts V,p)$.

\begin{Remark} %\label{rem:blowup}
The bound in~\eqref{lacase2} tends to infinity if $\be - \al \to 0$ unless $p=1$.
\end{Remark}


\medskip\noindent
{\bf Case ($\boldsymbol{ii}$).}
Using~\eqref{assumption=} to estimate~\eqref{beforecases} (as in the derivation of~\eqref{cba}),
we get
\begin{gather*}
%\MoveEqLeft
 \||I|^{-1} |\we a_k (t_0)| \|_{L^p (I)} + \big\|\we b'\big\|_{L^p(I)}
 \le C \big(\const\|1\|_{L^p(I)} + \|\we a'\|_{L^p (I)}\big), %\label{Lpestimate}
\end{gather*}
for a constant $C$ that depends only on $G \acts V$ and $p$; note that $B=B(G \acts V)$ by~\eqref{eq:constB}.
Thus, by~\eqref{amuab},
\begin{gather*}
 \|\ol a'\|_{L^p(I)} \le C \big(\const \|1\|_{L^p(I)} + \|\we a'\|_{L^p (I)}\big).
\end{gather*}
Let us set $A : = \max_{1 \le j \le n} \| a_j\|^{1/d_j}_{C^{d-1,1}([\al,\be])}$. Then
\begin{gather*}
 \const = \max_{1 \le j\le n} \big(\Lip_I\big(a_j^{(d-1)}\big)\big)^{1/d} |\we a_k (t_0)|^{(d-d_j)/d}
 \le \max_{1 \le j\le n} A^{d_j/d} A^{(d-d_j)/d} = A.
\end{gather*}
Consequently,
\begin{gather} \label{eq:laIA}
 \|\ol a'\|_{L^p(I)} \le C \big(A \|1\|_{L^p(I)} + \|\we a'\|_{L^p (I)}\big).
\end{gather}

 By Proposition~\ref{cover} (applied to $a$ instead of $b$ and
~\eqref{assumption=} instead of~\eqref{assumption2aG}),
 we can cover the set
 $(\al,\be) \setminus \{t \colon a(t) = 0\}$
by a countable family $\cI$ of open intervals $I$ on which~\eqref{eq:laIA} holds and such that
$\sum_{I \in \cI} |I| \le 2 (\be - \al)$.
Together with Lemma~\ref{lem:extend} we may conclude that
$\ol a$ is absolutely continuous on $(\al,\be)$ and satisfies
\begin{gather*}
 \|\ol a'\|_{L^p((\al,\be))} \le C \big(A \|1\|_{L^p((\al,\be))}
 + \|\we a'\|_{L^p ((\al,\be))}\big),
\end{gather*}
Using~\eqref{est} and the fact that $1 - 1/d_j < 1/p$ for all $j \le n$, we obtain
\begin{gather*}
 \|\ol a'\|_{L^p((\al,\be))}
 \\ \qquad
{} \le C \bigg(A (\be - \al)^{1/p}
+ \sum_{j=1}^n \max \big\{\big (\Lip_{(\al,\be)}\big( a_j^{(d_j-1)}\big)\big)^{1/d_j} (\be -\al)^{1-1/d_j}, \|a_j'\|^{1/d_j}_{L^\infty((\al,\be))}
 \big\} \bigg)
 \\ \qquad
{}\le C\, \max\big\{1, (\be-\al)^{1/p} \big\} \max_{1 \le j \le n} \|a_j\|^{1/d_j}_{C^{d-1,1}([\al, \be])},
\end{gather*}
where $C = C(G \acts V,p)$.
The proof of Theorem~\ref{main} is complete.
